\documentclass[utf8, a4paper, 14pt, russian, oneside]{book}

% Размер шрифта
\usepackage[14pt]{extsizes}

% Кодировка
\usepackage[T2A]{fontenc}
\usepackage[utf8]{inputenc}
\usepackage[main=russian, english]{babel}

\usepackage{fontspec}
\setmainfont{Times New Roman}
\usepackage{sectsty}
\sectionfont{\large}
\subsectionfont{\large}

% Параметры страницы
\usepackage[left=3cm, right=1cm, top=2cm, bottom=2cm]{geometry}
\pagestyle{plain}
\linespread{1.1} % Межстрочный интервал

% Пакеты для работы с математикой
\usepackage{amsmath}
\usepackage{amsfonts}
\usepackage{amssymb}

% Вставка изображений
\usepackage{graphicx}

% Пакет для работы с таблицами
\usepackage{tabularx}
\usepackage{booktabs}
\usepackage{longtable}

% Для больших множеств
\usepackage{mathtools}

% Для работы с рисунками
\usepackage{caption}

% Для создания графов в 3 блоке
\usepackage[all]{xy}

% Для специальных символов
\usepackage{textcomp}

% Для гиперссылок
\usepackage{hyperref}

% Для таблиц
\usepackage{multirow}

\usepackage{upgreek}

\usepackage{array}
\newcommand{\mysec}[1]{
{\centering\section*{\hyperlink{toc}{#1}}}
\addcontentsline{toc}{section}{#1}
}

\newcommand{\mysubsec}[1]{
{\centering\subsection*{\hyperlink{toc}{#1}}}
\addcontentsline{toc}{subsection}{#1}
}

\newcommand{\Db}{
    \ensuremath{
        D_{\text{в}}
    }
}

\newcommand{\xb}{
    \ensuremath{
        x_{\text{в}}
    }
}

\newcommand{\yb}{
    \ensuremath{
        y_{\text{в}}
    }
}

\newcommand{\zb}{
    \ensuremath{
        z_{\text{в}}
    }
}

\newcommand{\yx}{
    \ensuremath{
        y_x
    }
}

\newcommand{\zx}{
    \ensuremath{
        z_x
    }
}

\newcommand{\cov}{
    \ensuremath{
        \mathit{cov}_{\text{в}}
    }
}

\renewcommand{\r}{
    \ensuremath{
        \mathit{r}_{\text{в}}
    }
}

\newcommand{\rang}{
    \ensuremath{
        \mathit{rang}
    }
}

\newcommand{\rhob}{
    \ensuremath{
        \rho_{\text{в}}
    }
}

\newcommand{\der}[2]{
    \ensuremath{
        \frac{\partial #1}{\partial #2}
    }
}

% Команды для настройки содержания
\renewcommand\contentsname{\center{Содержание}} % Вместо оглавления пишется содержание
\addto{\captionsenglish}{\renewcommand{\bibname}{References}}
\begin{document}

\thispagestyle{empty}
~\vspace{-2cm}\setlength{\parindent}{0cm}
\begin{center}
	\includegraphics[scale=1.5]{../include/logo.png}\\[2pt]
	МИНОБРНАУКИ РОССИИ\\
	Федеральное государственное бюджетное образовательное учреждение\\
	высшего профессионального образования\\[5pt]
	\textbf{<<МИРЭА – Российский технологический университет>>}\\[5pt]
	\textbf{\large РТУ МИРЭА}\\[20pt]
	\hrule{}\mbox{}\\[1pt]
	\hrule{}\mbox{}\\[20pt]	
	Институт кибернетики \\ Кафедра <<Информационная безопасность>> (БК №252)\\[35pt]
	\textbf{Долгосрочное задание} \\
	по дисциплине: Математическая статистика
\end{center}
	\vspace{4in}
	Студент группы ККСО-01-19:  \qquad \qquad \qquad  \qquad     Колесников А.В.
\vspace{0.6in}
\begin{center}
Москва --- 2021
\end{center}
\newpage

\tableofcontents
\newpage

\mysec{Описание данных}

В качестве данных для анализа выберем оценки критиков, игроков, а такжу общую оценку вместе со всеми отзывами, для игр с сайта metacritic.com. Получим следующую выборку:
\begin{table}[h!]
    \centering
    \begin{tabular}{|c|c|c|c|}
        \hline
        № & $X$ & $Y$ & $Z$ \\ \hline
        1 & 89 & 66 & 89  \\ \hline
        2 & 83 & 80 & 95  \\ \hline
        3 & 82 & 94 & 85  \\ \hline
        4 & 81 & 96 & 93  \\ \hline
        5 & 73 & 77 & 87  \\ \hline
        6 & 95 & 79 & 91  \\ \hline
        7 & 84 & 99 & 87  \\ \hline
        8 & 93 & 80 & 84  \\ \hline
        9 & 84 & 92 & 81  \\ \hline
        10 & 85 & 72 & 93 \\ \hline
        11 & 89 & 76 & 96 \\ \hline
        12 & 76 & 86 & 80 \\ \hline
        13 & 81 & 69 & 83 \\ \hline
        14 & 69 & 68 & 99 \\ \hline
        15 & 91 & 85 & 89 \\ \hline
        16 & 68 & 77 & 87 \\ \hline
        17 & 96 & 77 & 97 \\ \hline
        18 & 76 & 78 & 95 \\ \hline
        19 & 77 & 87 & 90 \\ \hline
        20 & 97 & 93 & 98 \\ \hline
        21 & 79 & 78 & 99 \\ \hline
        22 & 93 & 97 & 98 \\ \hline
        23 & 93 & 96 & 81 \\ \hline
        24 & 91 & 83 & 83 \\ \hline
        25 & 92 & 85 & 92 \\ \hline
        26 & 67 & 78 & 81 \\ \hline
        27 & 89 & 90 & 82 \\ \hline
        28 & 71 & 83 & 86 \\ \hline
        29 & 71 & 79 & 88 \\ \hline
        30 & 98 & 84 & 91 \\ \hline
    \end{tabular} 
\end{table}

\newpage
\mysec{Корреляционный анализ}

\begin{gather*}
    % \xb = 83,77;\  \yb = 82,80;\  \zb = 89,33;
    \xb = \frac{89 + 83 + 82 + \ldots + 71 + 71 + 98}{30} = 83,77; \\
    \yb = \frac{66 + 80 + 94 + \ldots + 83 + 79 + 84}{30} = 82,80; \\
    \zb = \frac{89 + 95 + 85 + \ldots + 86 + 88 + 91}{30} = 89,33; \\
    \Db(x) = S^2(x) = \frac{(89 - 83,77)^2 + \ldots + (98 - 83,77)^2}{30} = 85,45; \\
    \Db(y) = S^2(y) = \frac{(66 - 82,80)^2 + \ldots + (84 - 82,80)^2}{30} = 76,23; \\
    \Db(z) = S^2(z) = \frac{(89 - 89,33)^2 + \ldots + (91 - 89,33)^2}{30} = 35,02; \\
    \sigma(x) = \sqrt{\Db(x)} = 9,24; \sigma(y) = \sqrt{\Db(y)} = 8,73; \sigma(z) = \sqrt{\Db(z)} = 5,92;
\end{gather*}

Найдем выборочные коэффициенты корреляции для каждой пары выборок:
\begin{gather*}
    \r(x,y) = \frac{\cov(x,y)}{\sigma(x)\sigma(y)} = 0,2881; \\
    \r(x, z) = 0,1720; \r(y, z) = -0,1781;
\end{gather*}

Далее найдем условные среднеквадратические отклонения:
\begin{gather*}
    S_{x/y} = \sqrt{S_x \cdot (1 - \r^2(x, y))} = 8,8517; \\
    S_{y/x} = 8,3606; S_{z/x} = 5,8297; S_{z/y} = 5,8234;
\end{gather*}

\newpage
\mysec{Частные коэффициенты корреляции}
Найдем частные коэффициенты корреляции по следующей формуле:
\begin{gather*}
    r_{xy/z} = \frac{
        \r(x,y) - \r(x, z)\r(y,z)
    }{
        \sqrt{1 - \r^2(x,z)}\sqrt{1 - \r^2(y, z)}
    } = 0,3288; \\
    r_{yz/x} = -0,3258; r_{xz/y} = 0,2370;
\end{gather*}

Сопоставим частные коэффициенты корреляции:
\begin{gather*}
    \r(x,y) < r_{xy/z} \Rightarrow \  \text{влияние усиливается}; \\
    \r(x,z) < r_{xz/y} \Rightarrow \  \text{влияние усиливается}; \\
    \r(y,z) > r_{yz/x} \Rightarrow \  \text{влияние ослабевает};
\end{gather*}

\mysec{Условное среднеквадратическое отклонение и множественный коэффициент корреляции}
Вычислим условное среднеквадратическое отклонение по следующей формуле:
\begin{gather*}
    S_{z/xy} = S_{z/x} \sqrt{1 - r^2_{yz/x}} = 5,5116
\end{gather*}
Найдем множественный коэффициент корреляции:
\begin{gather*}
    r_z = \sqrt{1 - \frac{S^2_{z/xy}}{S^2_z}} = 0,3642
\end{gather*}
Проверим полученный множественный коэффициент корреляции через парные коэффициенты корреляции:
\begin{gather*}
    r_z = \sqrt{
        \frac{
            \r^2(x, z) + \r^2(y, z) - 2\r(x,z)\r(y,z)\r(x,y)
            }{
                1 - \r^2(x, y)
        }
    } = 0,3642
\end{gather*}

\mysec{Значимость множественного коэффициента корреляции}
Проверим значимость множественного коэффициента корреляции $H_0: r^2_z = 0$ - множественный коэффициент корреляции незначим.

Вычислим статистику критерия:
\begin{gather*}
    F_{\text{набл}} = \frac{\frac{r_z^2}{2}}{\frac{1-r_z^2}{n-3}} = 2,0643; \\
    F_{\text{табл}}(\alpha=0,05; df_1=2;df_2=n-3) = 3,354; \\
    F_{\text{набл}} < F_{\text{табл}} \Rightarrow H_0\  \text{не отклоняется и}\  r^2_z \  \text{не значим}
\end{gather*}

\mysec{Уравнение регрессии}

Выпишем уравнение регрессии $z$ по $x, y$:
\begin{gather*}
    \overline{z}(x,y) - \overline{\zb} = b_{xz/y}(x - \overline{\xb}) + b_{yz/x}(y - \overline{\yb})
\end{gather*}

Найдем коэффициенты уравнения выше:
\begin{gather*}
    b_{xz/y} = r_{xz/y} \frac{S_{z/y}}{S_{x/y}} = 0,1559; \\ 
    b_{yz/x}=r_{yz/x}\frac{S_{z/x}}{S_{y/x}} = -0,2272; \\
    \overline{z}(x,y) - 89,33 = 0,1559(x - 83,77) - 0,2272(y - 82,80);\\
    \overline{z}(x,y) = 0,1559x - 0,2272y + 95,08
\end{gather*}

\mysec{Доверительные интервалы для коэффициентов множественной регрессии}
\begin{gather*}
    t_1=\frac{
        (b_{xz/y}-\beta_{xz/y})S_{x/y}\sqrt{n-3}
    }{
        S_{z/y}\sqrt{1-r_{zx/y}^2}
    } \\
    t_2=\frac{
        (b_{yz/x}-\beta_{yz/x})S_{y/x}\sqrt{n-3}
    }{
        S_{z/x}\sqrt{1-r_{zy/x}^2}
    }
\end{gather*}
Для построения доверительных интервалов нужно решить следующие
неравенства:
\begin{gather*}
\left|t_1=\frac{(b_{xz/y}-\beta_{xz/y})S_{x/y}\sqrt{n-3}}{S_{z/y}\sqrt{1-r_{zx/y}^2}}\right|<t_{\text{табл}} \\
\left|t_2=\frac{(b_{yz/x}-\beta_{yz/x})S_{y/x}\sqrt{n-3}}{S_{z/x}\sqrt{1-r_{zy/x}^2}}\right|<t_{\text{табл}}
\end{gather*}
, где $t_{\text{табл}}(\alpha = 0,05;df=n-3)=2,052$.

Решив неравенства получим следующие интервалы:
\begin{gather*}
    -0,0965 < \beta_{xz/y} < 0,4083 \\
    -0,4731 < \beta_{yz/x} < 0,0187
\end{gather*}

\mysec{Значимость частных коэффициентов корреляции}

Проверим значимость частных коэффициентов корреляции $H_0 : \rpriv = 0$.
\begin{gather*}
    t = \frac{
        \rpriv \sqrt{n-3}
    }{
        \sqrt{1 - \rpriv^2}
    } \\ 
    t_{xy/z} = 1,81; t_{xz/y} = -1,79; t_{yz/x} = 1,27;\\
    t_{\text{табл}}(\alpha=0,05; df = n-3) = 2,052; \\
    |t_{xy/z}| < t_{\text{табл}} \Rightarrow H_0 \  \text{принимается и коэффициент}\ r_{xy/z}\  \text{не значим}; \\
    |t_{xz/y}| < t_{\text{табл}} \Rightarrow H_0 \  \text{принимается и коэффициент}\ r_{xy/y}\  \text{не значим}; \\
    |t_{yz/x}| < t_{\text{табл}} \Rightarrow H_0 \  \text{принимается и коэффициент}\ r_{yz/x}\  \text{не значим};
\end{gather*}

\mysec{Построение доверительного интервала для частного коэффициента корреляции и множественного коэффициента корреляции}

Построим доверительный интервал для $r_{xy/z}$:
\begin{gather*}
    Z = \frac{1}{2}\ln\left(\frac{1+\rpriv}{1-\rpriv}\right) \approx N\left(
        \frac{1}{2} \ln \left( \frac{1+\rpriv}{1-\rpriv} \right); \frac{1}{\sqrt{n-4}}
    \right)
\end{gather*}

Тогда доверительный интервал для частного коэффициента корреляции находится при помощи неравенства:
\begin{gather*}
    \frac{1}{2} \ln \left(
        \frac{1+r_{xy/z}}{1-r_{xy/z}}
    \right)
    -
    \frac{1,96}{\sqrt{n-4}}
    <
    \frac{1}{2} \ln \left(
        \frac{1+\rho_{xy/z}}{1-\rho_{xy/z}}
    \right)
    <
    \frac{1}{2} \ln \left(
        \frac{1+r_{xy/z}}{1-r_{xy/z}}
    \right)
    +
    \frac{1,96}{\sqrt{n-4}}
\end{gather*}
Упростив выражение, получим следующее неравенство:
\begin{gather*}
    -0,0428
    <
    \frac{1}{2} \ln \left(
        \frac{1+\rho_{xy/z}}{1-\rho_{xy/z}}
    \right)
    <
    0,7259
\end{gather*}
В итоге получим, следующую систему:
\begin{gather*}
    \begin{cases}
        \rho_{xy/z} > -0,0428   \\
        \rho_{xy/z} < 0,62
    \end{cases}
\end{gather*}

Таким образом, доверительный интервал для частного коэффициента корреляции $r_{xy/z}: -0,04 < \rho_{xy/z} < 0,62$.
Построим доверительный интервал для множественного коэффициента корреляции $r_z$.
\begin{gather*}
    Z=\frac{1}{2}\ln\left(
        \frac{1+r_z}{1-r_z}
    \right) \approx N\left(
        \frac{1}{2}\ln\left(
            \frac{1+\rho_z}{1-\rho_z}
        \right);\frac{1}{\sqrt{n}}
    \right)
\end{gather*}
Значит доверительный интервал для множественного коэффициента корреляции находится с помощью следующего неравенства:
\begin{gather*}
    \frac{1}{2} \ln \left(
        \frac{1+r_{z}}{1-r_{z}}
    \right)
    -
    \frac{1,96}{\sqrt{n-4}}
    <
    \frac{1}{2} \ln \left(
        \frac{1+\rho_{z}}{1-\rho_{z}}
    \right)
    <
    \frac{1}{2} \ln \left(
        \frac{1+r_{z}}{1-r_{z}}
    \right)
    +
    \frac{1,96}{\sqrt{n-4}}
\end{gather*}
После подстановки значений получим следующее неравенство:
\begin{gather*}
    -0,003
    <
    \frac{1}{2} \ln \left(
        \frac{1+\rho_{z}}{1-\rho_{z}}
    \right)
    <
    0.7661
\end{gather*}
Решив неравенство получим:
\begin{gather*}
    \begin{cases}
        \rho_z > -0,003\\
        \rho_z < 0,645
    \end{cases}
\end{gather*}
Таким образом, доверительный интервал для множественного коэффициента корреляции $r_z: -0,003 < \rho_z < 0,645$.
 
\mysec{Корреляционная матрица. Множественный коэффициент корреляции.}
Построим корреляционную матрицу:
\begin{gather*}
    R = \begin{pmatrix}
        1 & \rhob(x,z) & \rhob(y,z) \\
        \rhob(x,z) & 1 & \rhob(x,y) \\
        \rhob(y,z) & \rhob(x,y) & 1
    \end{pmatrix}
    =
    \begin{pmatrix}
        1 & 0,1720 & -0,1781 \\
        0,1720 & 1 & 0,2881 \\
        -0,1781 & 0,2881 & 1
    \end{pmatrix}
\end{gather*}
Найдем множественный коэффициент корреляции:
\begin{align*}
    R^2_{z:xy} = 1 - \frac{|R|}{R_{11}} = 0,0861; \\
    r_z = \sqrt{R^2_{z:xy}} = 0,3642
\end{align*}

\mysec{Коэффициенты детерминации. Частные коэффициенты детерминации. Оценки коэффициентов линейной регрессии.}

Найдем коэффициенты детерминации:
\begin{gather*}
    R^2_{z:x} = 0,172^2 = 0,0296; R^2_{z:y} = -0.1781^2 = 0,0317; \\
    R^2_z = 0,3642^2 = 0,1326
\end{gather*}

Найдем частные коэффициенты детерминации:
\begin{gather*}
    R^2_{z:x\cdot y} = \frac{
        R_{z:xy^2} - R^2_{z:x}
    }{
        1 - R^2_{z:y}
    }
    \approx 0,0582; \\
    R_{z:y\cdot x}^2=\frac{
        R_{z:xy^2}-R_{z:y}^2
    }{
        1 - R_{z:y}^2
    }
    \approx 0,0562;
\end{gather*}

Сравним влияние факторов $x$ и $y$ на изменение $z$:
\begin{table}[h!]
    \centering 
    \begin{tabular}{|c|c|c|}
        \hline
        Факторы & Коэффициент детерминации & Частный коэффициент детерминации \\ \hline
        $x$ & 0,0296 & 0,0582 \\ \hline
        $y$ & 0,0317 & 0,0562 \\ \hline
        $z$ & 0,1326 &    \\ \hline
    \end{tabular}
\end{table}

Поскольку коэффициенты детерминации факторов $x$ и $y$ незначимо отличаются от коэффициента детерминации факторов $x, y$, то включение факторов $x$ и $y$ целесообразно.

Найдем оценки коэффициентов уравнения линейной регрессии:
\begin{gather*}
    Z = \begin{pmatrix}
        89 \\
        95 \\
        85 \\
        \vdots \\
        86 \\ 
        88 \\ 
        91
    \end{pmatrix};
    A = \begin{pmatrix}
        a_x \\
        a_y \\
        a_0
    \end{pmatrix};
    X = \begin{pmatrix}
        89 & 66 & 1 \\
        83 & 80 & 1 \\
        82 & 94 & 1 \\
        \vdots & \vdots & \vdots \\
        71 & 83 & 86 \\
        71 & 79 & 88 \\
        98 & 84 & 91
    \end{pmatrix}
\end{gather*}

Найдем оценки коэффициентов уравнения линейной регрессии $a_x, a_y, a_0$:
\begin{gather*}
    XA=Z \Leftrightarrow A = (X^TX)^{-1}X^TZ
\end{gather*}

Решив систему получим, что
\begin{gather*}
    A = \begin{pmatrix}
        0,1559 \\
        -0,1683 \\
        90,2034
    \end{pmatrix}
\end{gather*}

\mysec{Проверка значимости коэффициентов уравнения регрессии.}
Проверим значимость коэффициентов уравнения регрессии:
\begin{gather*}
    t_x = \frac{a_x}{S(a_x)};\\
    S(a_x)=\frac{\sigma_{\text{в}}(z)}{\sigma_{\text{в}}(x)}\sqrt{\frac{1-R_{z:xy}^2}{(1-R_{z:y}^2)(n-2-1)}} = 0,12; \\
    t_x = \frac{0,1559}{0,12} = 1,2784 \\
    t_y=\frac{a_y}{S(a_y)}; \\
    S(a_y)=\frac{\sigma_{\text{в}}(z)}{\sigma_{\text{в}}(y)}\sqrt{\frac{1-R_{z:xy}^2}{(1-R_{z:x}^2)(n-2-1)}} = 0,13; \\
    t_y = \frac{-0,1683}{0,13} = -1,3291; \\
    T_{\text{табл}}(df = n - 3) = 2,052; \\
    |t_x| < T_{\text{табл}} \Rightarrow H_0\  \text{принимается и } a_x \  \text{не значим}; \\
    |t_y| < T_{\text{табл}} \Rightarrow H_0\  \text{принимается и } a_y \  \text{не значим};
\end{gather*}
Переменные $x$ и $y$ не оказывают значительное влияние на $z$.

\mysec{Проверка значимости уравнения регрессии и дополнительных факторов.}
Проверим значимость уравнения регрессии $H_0:r_z^2=0$:
\begin{gather*}
    F_{\text{набл}}=\frac{r_z^2}{1-r_z^2}\frac{n-3}{2}=2,064; \\
    F_{\text{табл}}(\alpha = 0,05; df_1=m;df_2=n-m-1)=3,354; \\
    F_{\text{набл}} < F_{\text{табл}} \Rightarrow H_0\  \text{не отклоняется и уравнение регресии не значимо}.
\end{gather*}
Проверим значимость дополнительных факторов, включаемых в уравнение регрессии:
\begin{gather*}
    F_x=\frac{R_{z:xy}^2-\r^2(x,z)}{1-R_{z:xy}^2}\frac{n-3}{1} \approx 0,82; \\
    F_y=\frac{R_{z:xy}^2-\r^2(y,z)}{1-R_{z:xy}^2}\frac{n-3}{1} \approx 0,76; \\
    F_{\text{табл}}(\alpha = 0,05; df_1=1;df_2=n-m-1)=4,21; \\
    F_x < F_{\text{табл}} \Rightarrow H_0 \ \text{принимается и } x \ \text{не влияет на } z; \\
    F_x < F_{\text{табл}} \Rightarrow H_0 \ \text{принимается и } y \ \text{не влияет на } z;
\end{gather*}

\end{document}
